% PDF-derived LaTeX transcription. Original manuscript source was unavailable.
% Published PDF SHA-256: ebb9464aba4e7ff7e0b83df8e7d7d07a404f61c05bacc7019ac04bce9f037a75
% Extracted text SHA-256: dd343509ef1117cc62a6d58600e6578bc5e0ff46b47bcd36a81f5f273014d573
\documentclass[11pt]{article}
\usepackage[margin=1in]{geometry}
\usepackage{amsmath,amssymb}
\setlength{\emergencystretch}{3em}
\title{Discrete Constraints, Harmonic Modes, and Reconstruction Fidelity in Periodic Fluid and Magnetic Autoencoders}
\author{}
\date{}
\begin{document}
\maketitle
\noindent\textbf{PDF-derived LaTeX transcription.} This editable text follows the PDF extraction. Refer to the published PDF for figures, table layouts, and mathematical notation.

% Extracted chunk 1

Discrete Constraints, Harmonic Modes, and Reconstruction Fidelity
in Periodic Fluid and Magnetic Autoencoders

\section*{Abstract}

Divergence-free decoding restricts a reconstructed vector field but does not determine its
fidelity. We examine this distinction in a retrospective case study of periodic three-dimensional
fluid, magnetic, and scalar autoencoders with mixed continuous and finite-scalar-quantized la-
tents. We derive the consequences of the implemented unit-index centered-difference operators,
including the zero-mean obstruction of a periodic curl, the role of constant harmonic heads, and
the additional checkerboard null modes on even grids. We distinguish channel-relative energy ob-
jectives from total-energy and log-energy diagnostics. Historical training reports, contemporary
documentary gates, later reconstruction audits, and separate synthetic algebra tests are analyzed
without pooling their populations. In the later fluid audit, the validation-labeled set has com-
ponentwise MSE 2.912117 [U+0002] 10\ensuremath{-}6 and mean absolute discrete divergence 2.814595 [U+0002] 10\ensuremath{-}10, but
relative vorticity MSE 0.458089 fails the recorded reconstruction criterion. Every path in both
audit lists occurs in the retained training manifest, precluding an independently quarantined
validation interpretation. Multiscale magnetic comparisons exhibit modality tradeoffs rather
than a uniformly superior variant. A later joint candidate also retains substantial spectral and
rheological energy errors despite low divergence. The evidence supports a bounded conclusion:
construction can enforce a compatible discrete constraint, whereas mean recovery, amplitude,
spatial alignment, and cross-domain fidelity require distinct measurements and authenticated
evaluation populations.

\section*{1 Introduction and contributions}

Autoencoders for vector-valued scientific fields face two different requirements. Their outputs
should respect relevant structural constraints, and they should reconstruct the particular input
field accurately. These requirements are related but not interchangeable. A decoder that always
produces a discrete divergence-free velocity can still misplace vortices, suppress small-scale struc-
ture, or recover the wrong bulk flow. A magnetic decoder can satisfy the corresponding divergence
constraint while substantially misestimating a background field.
Periodic domains make this distinction especially transparent. Taking the curl of a periodic vec-
tor potential produces a divergence-free field under compatible differentiation, but its spatial mean
vanishes. A constant background velocity or magnetic field therefore needs a separate representa-
tional path. Moreover, a finite-difference implementation has a different null space from continuum
differentiation. These details matter when a structural identity is used to interpret small numerical
divergence as evidence of reconstruction quality.
Energy objectives introduce another distinction. The squared amplitude of a vector field can
be summarized by spatial means, spatial sums, separate channel contributions, or logarithmic
discrepancies. Although these quantities are related, their regularization and channel weighting
differ. An objective called volume-integrated energy need not equal either a physical energy integral
1

or the energy diagnostic used in an audit. Matching energy also does not fix the spatial arrangement
of that energy.
This paper studies these issues through a bounded retrospective analysis of periodic three-
dimensional fluid and magnetic autoencoders and their scalar branches. The contribution is not a
new divergence-free parameterization or a new quantization method. It is a joint account of which
properties follow from the captured construction, which reconstruction failures remain observable,
and how evaluation membership limits the inference.
Specifically, we provide a self-contained description of the implemented periodic operators, har-
monic and multiscale decoding paths, and hybrid quantizers; distinguish the principal reconstruc-
tion and energy metrics; and compare several levels of historical and later evidence without treating
them as interchangeable experiments. The resulting argument is deliberately narrower than a claim
of general-purpose physical representation: compatibility with one discrete divergence operator is
established separately from fidelity to the sampled fields.

\section*{2 Related work}

Richter-Powell and colleagues establish neural constructions that satisfy divergence and continuity
constraints by design [ 1]. Their periodic continuum theorem explicitly allows a missing constant
vector field. This is directly relevant to the bulk-flow and magnetic-background obstruction con-
sidered here. We refer to the 2022 NeurIPS paper through its arXiv version 3. Its continuum
expressiveness result does not establish universality of a finite-resolution, quantized decoder with
a centered-difference curl.
Finite scalar quantization bounds individual latent coordinates and maps them to fixed scalar
levels, thereby defining an implicit product codebook [ 2]. Straight-through training supplies an
approximate gradient through quantization. We use the 2023 arXiv version 2 preprint for this
methodological context. The mixed continuous/SplitFSQ quantizers in the present case study are
a particular implementation with continuous coordinates, scalar-level assignments, and an addi-
tional utilization proxy. Neither the cited FSQ construction nor observed marginal bin occupancy
demonstrates that every joint code combination occurs in this implementation. No external bench-
mark values from either reference are used.

\section*{3 Scientific methods and validation protocols}


\section*{3.1 Retrospective scope and evidence populations}

The analysis uses captured source text, training reports, audit reports, manifests, and independently
generated finite algebra diagnostics. No retrieved code was executed, no checkpoint or pickle
was loaded, and no original raw field array was loaded during this reanalysis. There was no
additional neural training, neural inference, or physical simulation. The new calculations concern
saved-record consistency and separately implemented synthetic algebra, not a replay of the original
computational environment.
We distinguish historical training-final losses; a contemporary documentary gate with train- and
validation-labeled lists; a later actual reconstruction audit; a separate stress audit; per-domain
magnetic and scalar comparisons; a later joint candidate and rheology audit; and new synthetic
2

Table 1: Evidence populations and their interpretation. Counts are not pooled.
Evidence Recorded scope Interpretation
Fluid training-final Training reduc-
tion
Optimization record, not audit accuracy
Contemporary fluid gate 100 + 100 paths Earlier documentary contract
Later fluid reality audit 100 + 100 paths Actual saved reconstruction audit
Later fluid stress audit 64 cases Perturbation and failure diagnostics
Multiscale comparisons 64 per domain Within-source Hartmann/Arrhenius com-
parison
Later joint core audit 64 per domain Separate candidate and training history
Rheology audit 128 cases Fluid head on trajectory channels
Synthetic algebra 40 cases and two
grids
Finite arithmetic, not neural validation
operator tests. Table 1 summarizes the populations most important to interpretation. Counts are
numbers of recorded cases or paths, not authenticated independent physical trajectories.
All 100 later fluid audit train paths and all 100 validation-labeled paths occur in the retained
4096-store training manifest. The two audit lists are path-disjoint, but each spans 83 parent
directories and they share 22 parent directories. Thus the labels do not demonstrate independently
quarantined validation. Directory grouping is not independent authentication of physical identity,
and membership in earlier warm-start training is not fully reconstructed.
For the multiscale comparisons, the captured loader selects the first N sorted Hartmann and
Arrhenius paths. Training reports list 5000 samples per domain, while the audit selects the first 64
per domain; the source contains no exclusion split. Exact historical execution and array identities
remain incomplete. We therefore use the term audit set, not held-out set.

\section*{3.2 Periodic discrete operators}

Let G = ZN1 [U+0002] ZN2 [U+0002] ZN3 be a periodic lattice with V = N1N2N3 sites. A scalar field is f : G ! R,
and a vector field is u : G ! R3. Indices wrap modulo the corresponding grid length. For coordinate
unit vectors ej, define
Djf (n) = f (n + ej) [U+0000] f (n [U+0000] ej)
2 , j = 1, 2, 3. (1)
These derivatives use unit grid indices, without physical spacing factors. Define discrete divergence
D and curl C by
Du = D1u1 + D2u2 + D3u3, (2)
CA = (D2A3 [U+0000] D3A2, D 3A1 [U+0000] D1A3, D 1A2 [U+0000] D2A1). (3)
The shift operators commute, so DiDj = DjDi and DCA = 0 in exact arithmetic. With ¯f =
V \ensuremath{-}1 P
n\ensuremath{\in}G f (n), periodic telescoping also gives Djf = 0 and CA = 0 . Floating-point implemen-
tations leave small residuals. Reported divergence values are diagnostics of this operator, not
independently calibrated SI measurements.
A basic vector decoder therefore has the form
bu = \ensuremath{\alpha}CA + h, h 2 R3, (4)
3

where \ensuremath{\alpha} is spatially uniform within a sample and h is broadcast over the grid. The constant
harmonic head supplies a mean that a periodic curl cannot represent. For a pure-curl prediction,
Jensen’s inequality gives
1
3V
X
n
kbu(n) [U+0000] u(n)k2 [U+0015] 1
3 k¯uk2. (5)
More generally, writing u\ensuremath{\prime} = u [U+0000] ¯u and bu\ensuremath{\prime} = bu [U+0000] bu yields the exact decomposition
M (bu, u) = M (bu\ensuremath{\prime}, u\ensuremath{\prime}) + 1
3 kbu [U+0000] ¯uk2. (6)
Here M is the componentwise MSE defined below. Replacing the predicted mean with the target
mean is an oracle diagnostic of this decomposition, not a deployable reconstruction procedure.
On an even lattice, centered differentiation has additional null modes. Its Fourier symbol is
i sin(2\ensuremath{\pi}kj/Nj), which vanishes at both zero and the Nyquist frequency. At nonzero frequency
vectors whose components are each zero or Nyquist, the entire curl symbol vanishes. For example,
q(n) = (( [U+0000]1)n1, 0, 0) has zero discrete divergence and zero mean but cannot be represented as
a periodic centered-difference curl plus a constant. A constant harmonic head repairs the mean
obstruction, not every discrete null-space obstruction.

\section*{3.3 Gains, gates, and additional solenoidal paths}

A spatially uniform scalar a and constant vector b preserve divergence under the isotropic affine
map u 7! au + b. Unequal channel gains generally do not: if G = diag(g1, g2, g3), then
D(Gu) =
X
j
gjDjuj, (7)
which need not vanish when P
j Djuj = 0 . For a spatial scalar gate s(n), the exact discrete
expression is
D(su)(n) = 1
2
X
j
[s(n + ej)uj(n + ej) [U+0000] s(n [U+0000] ej)uj(n [U+0000] ej)] . (8)
There is no general cancellation. In contrast, applying a gate to the potential before differentiation
gives C(sA), whose matching divergence still vanishes. Componentwise scaling of the potential also
remains inside the curl and is distinct from channel scaling of the decoded vector.
The captured single-scale source additionally permits profiles p1(n2, n3), p2(n1, n3), and p3(n1, n2).
Each component is constant along its own derivative direction, so their combined divergence is zero.
These profiles and optional potential-based refinement preserve the structural identity. Their pres-
ence in captured source does not authenticate their exact historical activation in every checkpoint.

\section*{3.4 Hybrid multiscale architecture}

The single-scale reference maps 1283 inputs to a 163 latent lattice with 48 channels, of which 16 use
the continuous path. The captured multiscale architecture first applies a three-dimensional Haar
transform, followed by convolutions and strided downsampling. Its latent lanes have spatial sizes
and channel counts 83/32, 163/48, and 323/24. They respectively contain 16, 16, and 12 continuous
channels; the remaining coordinates are scalar-quantized independently within each lane.
4

For a continuous latent coordinate z, the implemented map is qc(z) = 2 tanh( z/2). For a discrete
coordinate, let t = tanh(z) and
L = f[U+0000]1, [U+0000]2/3, [U+0000]1/3, 0, 1/3, 2/3, 1g, Q (t) = argmin
\ensuremath{\ell}\ensuremath{\in}L
jt [U+0000] \ensuremath{\ell}j. (9)
The forward value is Q(t), implemented as t+stopgrad(Q(t)[U+0000]t) for straight-through differentiation.
Discrete channels are partitioned into four groups with the same seven levels. A soft assignment
proportional to exp[[U+0000](t[U+0000]\ensuremath{\ell})2/0.05] defines a marginal normalized-entropy proxy. This proxy and the
recorded quantization discrepancy are implementation-specific diagnostics. They do not measure
exact retained information or a stored-bit compression rate.
Writing quantized lanes as z8, z16, z32, the multiscale vector heads have the forms
bu = CAu(z16) + hu(z8) + 0.25C [s(z32)As(z32)] , (10)
bB = CAB(z8, z16, z32) + hB(z8). (11)
The scalar decoder directly fuses all lanes and has five outputs in the audited configuration. The
radiation decoder is a direct six-channel tensor head. Its existence does not demonstrate trained ra-
diation reconstruction or physical tensor validity. Fluid and magnetic inputs are routed separately;
the Arrhenius branch conditions on temperature and two reactant concentrations and predicts those
fields together with a third concentration and reaction rate.

\section*{3.5 Metrics and energy objectives}

For C-channel fields x, bx : G ! RC, define
M (bx, x) = 1
CV
X
n,c
(bxc(n) [U+0000] xc(n))2. (12)
This is the raw reconstruction metric. The fluid audit also reports
R2 = kbu [U+0000] uk2
kuk2 + 10\ensuremath{-}12 , (13)
R\ensuremath{\omega} = M (Cbu, Cu)
M (Cu, 0) + \ensuremath{\varepsilon}\ensuremath{\omega}
. (14)
Norms include all sites and components. The later fluid evaluator uses \ensuremath{\varepsilon}\ensuremath{\omega} = 10 \ensuremath{-}30; the multiscale
reported vorticity loss uses 10\ensuremath{-}10. These are distinct regularization conventions. R\ensuremath{\omega} is a relative
MSE, not a relative RMS error, and cohort means cannot generally be converted by taking a single
square root.
Let Sc(u) = P
n uc(n)2 and K(u) = 1
2
P
c Sc(u). The implemented channel-sum objective is
LchE = 1
3
3X
c=1
jSc(bu) [U+0000] Sc(u)j
Sc(u) + 10\ensuremath{-}8 . (15)
It has no explicit density or cell-volume factor. It averages channel-relative discrepancies rather
than taking the relative discrepancy of total vector energy. Separate audit diagnostics are
RK = jK(bu) [U+0000] K(u)j
jK(u)j + 10\ensuremath{-}12 , (16)
AK =

% Extracted chunk 3

log max(K(bu), 10\ensuremath{-}30) [U+0000] log max(K(u), 10\ensuremath{-}30)

% Extracted chunk 5

. (17)
5

A captured training loss instead squares the difference between logarithms of spatial-mean energies
with additive 10\ensuremath{-}12 regularization. Multiscale training includes absolute differences of logarithms
of total vector squared amplitudes with additive 10\ensuremath{-}8 regularization. These definitions are not
automatically equivalent to AK.
For constant density \ensuremath{\rho} and uniform cell volume \ensuremath{\Delta}V , a physical kinetic energy would be \ensuremath{\rho}\ensuremath{\Delta}V K(u).
Common factors cancel from an unregularized relative ratio and from a log-energy difference. With
regularization, cancellation requires scaling the regularizer consistently. Likewise, converting sums
to means requires scaling additive regularizers. There is no inevitable relative-error growth pro-
portional to voxel count. Under a uniform gain a, K(au) = a2K(u), while divergence can remain
zero.
The radial spectral discrepancy uses an unnormalized forward real-input discrete Fourier trans-
form. Power is averaged over channels and then over the stored half-spectrum coefficients in each
of 12 radial frequency bins, without conjugate-multiplicity weights. This defines a bin-averaged
diagnostic rather than a shell-integrated physical energy spectrum. With channel-averaged power
Px(k) and bins Bb, it is
Lrad = 1
jIj
X
b\ensuremath{\in}I
[log(1 + hPbxiBb) [U+0000] log(1 + hPxiBb)]2 , (18)
where I contains nonempty bins. Frequencies are in index units; the outer edge uses the captured
strict upper-bound convention. This statistic discards Fourier phase. The local vorticity diagnos-
tic instead compares centered-difference curls after circular 33 averaging using an absolute-error
reduction. Energy, radial spectra, voxelwise vorticity, and local smoothed vorticity therefore test
different properties.

\section*{3.6 Preparation, arithmetic, and audit decisions}

The captured fluid reader selects velocity or trajectory channels, takes the first frame when ap-
plicable, converts to float32, replaces nonfinite values by zero, and resizes to 1283 by trilinear
interpolation when needed. The later fluid audit normalizes inputs with a constant channel-mean
vector and one shared standard deviation, then denormalizes outputs. The shared scale preserves
the discrete divergence property. The multiscale reader converts channel-last Hartmann arrays to
channel-first float tensors and resizes when required. These preparation steps can affect the target
field and should not be confused with physical calibration.
The retained multiscale training source uses AdamW and supports CUDA bfloat16 autocast
with gradient scaling. Captured audit functions use evaluation mode without an explicit autocast
context. Exact historical package versions and runtime bindings are not established by source text
alone.
The later fluid vorticity gate requires mean R\ensuremath{\omega} [U+0014] 0.25 and its casewise 95th percentile [U+0014] 0.60.
Its log-energy limits are 0.10 for the mean and 0.25 for the 95th percentile. Divergence limits are
10\ensuremath{-}7 for mean absolute divergence and 10\ensuremath{-}6 for the 95th percentile of casewise maxima. These are
recorded diagnostic contracts, not externally validated physical tolerances.
Separate new float64 algebra calculations use NumPy 2.3.5 and a predeclared arithmetic tolerance
of 10\ensuremath{-}10. The random tests comprise 20 trials on each of two periodic grids, 163 and 243, for 40 cases
in total. Separate analytic counterexamples are evaluated on both grid sizes. These calculations
check operator cancellation, mean-error decomposition, gain behavior, and energy scaling; they test
the stated algebra locally, not the learned models.
6

Table 2: Later fluid reconstruction audit. Except for case counts, entries are case means unless
marked p95 (the casewise 95th percentile). Both labeled lists occur in the retained training manifest
and do not constitute independently quarantined validation.
Metric Train-labeled Validation-labeled
Cases 100 100
Componentwise MSE 2.730878 [U+0002] 10\ensuremath{-}6 2.912117 [U+0002] 10\ensuremath{-}6
Relative L2 0.294821 0.296723
Relative energy error 0.052315 0.042931
Absolute log-energy error 0.033557 0.030256
Relative vorticity MSE 0.474494 0.458089
Relative vorticity MSE, p95 0.549369 0.613239
Radial spectral discrepancy 0.012003 0.011863
Mean absolute divergence 2.786640 [U+0002] 10\ensuremath{-}10 2.814595 [U+0002] 10\ensuremath{-}10
Target mean absolute divergence 6.371008 [U+0002] 10\ensuremath{-}5 6.398082 [U+0002] 10\ensuremath{-}5

\section*{4 Experimental results}


\section*{4.1 Training-final values and documentary gates}

The selected energy-repaired fluid reference has a historical training-final reconstruction loss of
1.667281 [U+0002] 10\ensuremath{-}6 and relative L2 value 0.129984. These are training reductions, not the later
audit measurements. A second learning-rate continuation records reconstruction 1.688426 [U+0002] 10\ensuremath{-}6
and relative L2 0.148683. Their similar raw MSEs do not establish equivalent energy or relative
reconstruction behavior.
The contemporary documentary gate records a pass on 100 train- and 100 validation-labeled
paths. For its validation-labeled set it reports MSE 2.912083 [U+0002] 10\ensuremath{-}6, relative L2 0.296719, rela-
tive energy error 0.042925, energy-error 95th percentile 0.047140, and mean absolute divergence
2.813930 [U+0002] 10\ensuremath{-}10. These are retained as documentary values, not substituted for the later actual
JSON values in Table 2.

\section*{4.2 Low divergence does not resolve vorticity error}

The later reality audit fails its voxelwise-vorticity check. On the validation-labeled set, both mean
relative vorticity MSE 0.458089 and its 95th percentile 0.613239 exceed the respective limits. Low
divergence, low absolute log-energy error, and a small radial spectral discrepancy do not remove
this failure. The target itself has nonzero divergence under the same stencil, so exact solenoidality
of the prediction also need not reproduce every target component.
The separate 64-case stress audit records MSE 2.397006 [U+0002] 10\ensuremath{-}6, mean absolute divergence
2.697498 [U+0002] 10\ensuremath{-}10, and absolute log-energy error 0.014545. Nevertheless, relative vorticity MSE
is 0.415995. The mean fitted spectral slope is [U+0000]5.374599 for predictions versus [U+0000]4.034473 for tar-
gets, with mean absolute slope error 1.340127. The stress gate fails spectral-slope, vorticity, and
phase-interface checks. These fitted slopes are diagnostics on the captured fields, not evidence that
either field realizes an independently validated inertial-range law.
The phase-interface diagnostic gives a gradient-steepness ratio of 32.7161, but it applies a velocity-
only checkpoint to a scalar \ensuremath{\varphi} replicated into three channels. It is a proxy failure, not an indepen-
7

Table 3: Multiscale audit comparisons: 64 cases per domain. Fluid vorticity is relative MSE;
reconstruction columns are componentwise MSE.
Model Fluid MSE Fluid vort. Magnetic MSE Scalar MSE
Pre-routing baseline 3.8785 [U+0002] 10\ensuremath{-}4 0.02575 2.01278 3 .9552 [U+0002] 10\ensuremath{-}4
Magnetic-input routing 4.3062 [U+0002] 10\ensuremath{-}4 0.04449 0.70913 2 .7665 [U+0002] 10\ensuremath{-}3
Harmonic anchor 2.6606 [U+0002] 10\ensuremath{-}4 0.02431 0.33804 6 .1233 [U+0002] 10\ensuremath{-}4
Harmonic energy-weighted 3.3547 [U+0002] 10\ensuremath{-}4 0.05506 0.54570 7 .6626 [U+0002] 10\ensuremath{-}4
Sharp harmonic balance 2.8701 [U+0002] 10\ensuremath{-}4 0.02255 0.51936 4 .7869 [U+0002] 10\ensuremath{-}4
dently validated phase-field reconstruction experiment. The important conclusion is that small
raw MSE and a compatible divergence constraint did not certify local reconstruction in this stress
population.

\section*{4.3 Gain alternatives and evaluator ambiguity}

The per-channel post-curl gain training report records divergence 1.135488 [U+0002] 10\ensuremath{-}6, larger than
2.257618 [U+0002] 10\ensuremath{-}10 in the selected reference training report. Its contemporary Markdown audit, how-
ever, records validation mean absolute divergence 2.790532[U+0002]10\ensuremath{-}10. These records do not constitute
a matched empirical gain ablation. The captured generic evaluator constructs a default model and
permits partial checkpoint loading; the exact historical gain-path evaluation is unresolved. Partial
loading is a source-level possibility, not an established explanation for the discrepancy.
The isotropic-gain documentary audit preserves low divergence, 3.079795 [U+0002] 10\ensuremath{-}10, while its val-
idation relative energy error is 0.089741 and its gate fails on energy. A separate scratch-training
document records final MSE 1.0982 [U+0002] 10\ensuremath{-}4 and radial spectral discrepancy 2.099282. These partic-
ular alternatives do not establish universal inadequacy of isotropic gains or scratch training. The
algebraic distinction between isotropic and anisotropic post-curl gains remains valid independently
of their historical evaluator binding.

\section*{4.4 Magnetic harmonic anchoring and modality tradeoffs}

Table 3 reports the saved multiscale comparisons using descriptive model names. The sequence
includes correction of magnetic input routing and introduction of direct supervision of the harmonic
magnetic head against the spatial target mean. These are combined historical modifications, not
an isolated causal intervention.
The harmonic-anchor variant has lower magnetic MSE than the sharp harmonic-balance variant,
0.33804 versus 0.51936. The latter instead improves relative fluid-vorticity MSE and scalar MSE
within this comparison. Its magnetic mean absolute divergence is 5.284846 [U+0002] 10\ensuremath{-}8, but this struc-
tural residual is not a ranking of magnetic reconstruction accuracy. The multiscale vorticity value
0.02255 is a relative MSE, not a percentage RMS error.
The saved magnetic harmonic MSE is nearly equal to full magnetic MSE for the harmonic-anchor
variants. For example, the harmonic-anchor report gives 0.3380369043 for the harmonic error and
0.3380369041 for full reconstruction. This is consistent with a dominant mean contribution at the
recorded precision, but the absent original arrays prevent a complete spatial-error decomposition.
Training-final magnetic reconstruction for that variant is 1.062824, distinct from its audit value.
8

The separation between training reductions and audit summaries remains necessary even within
one named variant.

\section*{4.5 Later joint candidate}

The later joint candidate is evaluated from a recorded resume artifact; a directory label does
not authenticate cumulative optimizer steps. In its separate core audit, fluid MSE is 0.029359,
relative vorticity MSE is 1.332971, and radial spectral discrepancy is 14.7002, while mean absolute
divergence remains 1.405821 [U+0002] 10\ensuremath{-}8. Its magnetic MSE is 0.067119, but magnetic radial spectral
discrepancy is 9.678850. Better performance on one amplitude-sensitive summary does not establish
uniformly better field reconstruction.
The separate 128-case rheology audit applies the fluid head to trajectory channels corresponding
to velocity. It records mean absolute log-energy error 2.88487, relative vorticity MSE 1.031074, and
mean absolute divergence 1.758416 [U+0002] 10\ensuremath{-}9. The population contains 55 Carreau–Yasuda-labeled
and 73 power-law-labeled cases, without an authenticated independent-trajectory count. These
adverse energy and vorticity results remain part of the assessment despite the small divergence
residual.

\section*{4.6 Independent finite algebra diagnostics}

The separate synthetic calculations confirm the expected cancellations and counterexamples at
the declared arithmetic tolerance. They support zero divergence of a compatible curl, zero mean
of a periodic curl, the mean-error decomposition, and quadratic energy scaling under uniform
gain. They also demonstrate nonzero divergence after anisotropic post-curl gains and spatial scalar
gates, and the centered-difference checkerboard obstruction. These are finite local demonstrations
of standard identities and failure mechanisms, not original neural inference or physical validation.
Exact inputs, recipes, residuals, and saved-record checking details are retained in Supplementary
record S1.

\section*{5 Discussion and limitations}


\section*{5.1 What is guaranteed, and what remains to be learned}

The compatible curl/divergence pair supplies a meaningful construction-level property: in exact
arithmetic, all decoded curl contributions have zero matching discrete divergence. Constant bulk
terms and the additional explicitly solenoidal paths preserve it. This property is useful because it
does not depend on optimization successfully minimizing a divergence penalty. It remains limited
to the specified operator, boundary convention, and transformations applied after decoding.
Harmonic heads address a separate representational necessity. A nonzero spatial mean cannot
pass through a periodic curl, and energy matching cannot recover a missing mean if the architecture
forbids it. Equation 6 makes the resulting error floor explicit. Nevertheless, adding a constant does
not make the centered-difference decoder complete on even grids: Nyquist null modes remain. Con-
tinuum representation results and discrete decoder expressiveness must therefore be kept separate.
The measurements demonstrate why structural feasibility should not serve as a proxy for fidelity.
Energy is a scalar summary; radial spectral losses further discard spatial phase. Neither controls the
location or sign alignment of individual eddies. The later fluid audit provides a concrete example
9

in which divergence and energy diagnostics are favorable while voxelwise vorticity reconstruction
fails. The stress and later joint-candidate results extend this distinction under different populations
and contracts rather than replicating a single common accuracy experiment.

\section*{5.2 Energy normalization and multimodal interpretation}

Equation 15 places equal outer weight on channel-relative discrepancies. Consequently, a weak
transverse component may matter strongly to that objective even when a dominant longitudinal
or background component controls total energy. Conversely, errors with opposite signs across
channels can partially cancel in total-energy differences while remaining visible to the channel-
relative objective. Log-energy diagnostics are different again, especially near a regularization floor.
Density, cell volume, spatial sums, and spatial means must be tracked consistently, but a large
grid does not by itself imply a large relative energy error. The evidence does not support an
inevitable voxel-count amplification mechanism. The objective’s scientific interpretation instead
depends on its denominators, channel aggregation, field normalization, and treatment of small
target energies.
The multiscale comparisons also resist a single overall ranking. The harmonic-anchor and sharp
harmonic-balance variants trade magnetic reconstruction against selected fluid and scalar metrics.
Routing, supervision, loss weights, and training histories are not isolated. Thus the comparison
describes observed behavior within the audit source, not a causal estimate of the benefit of any one
architectural change. Scalar channels also represent different quantities; an all-channel MSE is a
numerical summary of the captured preparation, not an automatically dimensionally commensurate
physical error.

\section*{5.3 Evaluation scope and reproducibility limits}

The central generalization limit is explicit training-manifest overlap. Path-disjoint audit lists can
still contain training members, shared parent groups, or related states of the same physical system.
Here every later fluid audit path occurs in the retained training manifest. The multiscale selection
code likewise does not establish an exclusion split. Directory names cannot independently resolve
physical identity, and prior warm-start membership remains incompletely reconstructed. We there-
fore provide no pooled success rate, confidence interval based on assumed independent trajectories,
or independently quarantined generalization estimate.
The captured evaluator and model sources are sufficient to describe intended operations, but they
do not fully authenticate historical runtime bindings. In particular, the inconsistent per-channel-
gain divergence records cannot be resolved into a matched ablation without the relevant artifact
and loader trace. The later joint candidate’s checkpoint provenance likewise does not establish its
cumulative optimizer duration.
Saved-record consistency checks support many transcriptions and aggregate calculations, but
small differences between separately reduced scalar summaries remain. They are retained without
changing scientific thresholds. Original scalar arrays and intermediate reductions were not replayed.
The finite synthetic tests establish arithmetic behavior only over their declared cases; they do not
supply missing physical calibration, solver validation, radiation validation, or learned reconstruction
evidence.
10


\section*{6 Conclusion}

Periodic fluid and magnetic autoencoders can enforce a compatible discrete divergence constraint
through curl-based decoding, but this construction neither guarantees a correct spatial mean nor
certifies reconstruction fidelity. Constant harmonic heads restore the mean path, while even-grid
centered-difference null modes delimit the resulting representation. Uniform isotropic scaling pre-
serves the divergence identity; unequal post-curl channel gains and spatial scalar gates generally
do not.
The retrospective measurements reinforce these distinctions. The later fluid audit fails its
voxelwise-vorticity criterion, and the separate stress audit fails both vorticity and spectral-slope
checks, despite small raw MSE, low divergence, and favorable energy summaries. Multiscale mag-
netic comparisons show modality tradeoffs, and the later joint candidate retains substantial spec-
tral and rheological energy errors. Training-manifest overlap and incomplete historical runtime
authentication restrict these results to bounded audit comparisons. A scientifically interpretable
evaluation must therefore report construction-level constraints, harmonic mean recovery, energy
normalization, spatial reconstruction diagnostics, and population provenance separately.

\section*{7 Data and Code A vailability}

Supplementary record S1 is a private export containing original source, report, and manifest cap-
tures; exact local checking inputs, scripts, and results; and the two primary reference PDFs. The
source captures include model, loss, quantizer, loader, training, and evaluation code as text. Sepa-
rate scripts implement the local saved-record checks and synthetic-algebra calculations. The export
is not a public archive, and no public or reviewer access mechanism or permission to redistribute
the original project artifacts is documented.
The original neural checkpoints and raw field arrays were not loaded in this reanalysis. Complete
reproduction of the original training and inference environment is not established by the private text
captures. The separate local synthetic-algebra and saved-record calculations have retained inputs
and recipes in Supplementary record S1; their reproducibility scope is those finite calculations, not
the original neural experiments. The cited primary papers are publicly available through their
arXiv records.

\section*{References}

[1] J. Richter-Powell, Y. Lipman, and R. T. Q. Chen. Neural Conservation Laws: A Divergence-Free
Perspective. Advances in Neural Information Processing Systems 35 (2022). arXiv:2210.01741v3.
https://arxiv.org/abs/2210.01741.
[2] F. Mentzer, D. Minnen, E. Agustsson, and M. Tschannen. Finite Scalar Quantization: VQ-V AE
Made Simple. arXiv:2309.15505v2 (2023). https://arxiv.org/abs/2309.15505.
11

\end{document}
